\section{An exact terminal test and the two-squared-axis certificate}
\label{subsec:terminal-grid-certificate}

The next criterion is a linear feasibility statement for a prescribed
signed difference. Its use here is to exclude homometric divisor ideals.
Feasibility alone does not construct such ideals.

Let \(K=\prod_i p_i^{a_i}\) be odd, with physical order grid
\(\mathcal G=\prod_i\{0,\ldots,a_i\}\). For a nonzero radial
\(d:\mathbb Z_K\to\{-1,0,1\}\), let \(U\) be the upward closure
of its support and \(Z\) its zero set of Fourier conductors. Index
conductors by their coordinate kernel exponents \(h\in\mathcal G\).
We call \(d\) \emph{terminally feasible} if a real radial function
\(g\) satisfies
\begin{equation}
g(v)=|d(v)|\quad(v\in U),\qquad
\widehat g(h)=0\quad(h\notin Z).
\label{eq:terminal-feasibility}
\end{equation}
No binary or monotonicity condition is placed on \(g\).

\begin{lemma}[Integer incidence test]
\label{lem:terminal-incidence}
Extend functions by zero beyond the top of the grid and put
\[
b_v=\sum_{\substack{\epsilon\in\{0,1\}^n\\v+\epsilon\in\mathcal G}}
       (-1)^{|\epsilon|}|d(v+\epsilon)|,
\qquad
M_{vh}=\prod_i\bigl(\mathbf1_{h_i=v_i}-\mathbf1_{h_i=v_i-1}\bigr).
\]
Then \eqref{eq:terminal-feasibility} is feasible if and only if the
real linear system
\begin{equation}
\frac{b_v}{K}\prod_i p_i^{v_i}
       =\sum_{h\in Z}M_{vh}t_h\qquad(v\in U)
\label{eq:terminal-incidence-system}
\end{equation}
is solvable. In particular, every integer left-null vector \(w\)
of \(M\) imposes the necessary identity
\[
\sum_{v\in U}w_vb_v\prod_i p_i^{v_i}=0.
\]
\end{lemma}

\begin{proof}
On one axis, the inverse Fourier column indexed by kernel exponent
\(h\) is the character sum \(I_h(v)\) over the characters of order
\(p^{a-h}\) at a point of physical order \(p^v\). It has values
\(\varphi(p^{a-h})\) for \(v\le h\),
\(-p^{a-h-1}\) for \(v=h+1\), and zero for \(v>h+1\).
For \(h=a\) it is identically one. Direct subtraction, including
the top boundary with zero extension, gives
\[
I_h(v)-I_h(v+1)
=p^{a-v}\bigl(\mathbf1_{h=v}-\mathbf1_{h=v-1}\bigr).
\]
Write the inverse transform as
\(g(v)=\sum_{h\in Z}t_h\prod_i I_{h_i}(v_i)\), so
\(t_h=\widehat g(h)/K\). Taking the product of the forward
differences gives \eqref{eq:terminal-incidence-system} wherever
\(g=|d|\) on all upper neighbors. This proves necessity because
\(U\) is upward closed. Conversely, a solution defines a radial
\(g\) with the required Fourier support. Its adjacent differences
agree with those of \(|d|\) on \(U\). Recovering values from the
top downward in this triangular difference system proves equality on
\(U\). Multiplying the system by a left-null vector and then by
\(K\) gives the final identity.
\end{proof}

\begin{lemma}[Why terminal exclusion implies rigidity]
\label{lem:terminal-implies-rigidity}
If every nonzero equal-population difference of two divisor downsets
on \(K\) is terminally infeasible, then \(2^eK\) is rigid for
every \(e\ge0\).
\end{lemma}

\begin{proof}
For a purported pair take the least disputed dyadic size \(L\) and
its signed indicator \(d\). Lemma~\ref{lem:common-dyadic-fibers}
gives equal positive and negative populations. Its sign supports are
mutually incomparable. Each is order-convex: between two same-sign
size-\(L\) disputes the empty side stays empty and the occupied
size, nonincreasing between equal endpoints, stays \(L\).
Let \(D\) be the downward closure of both sign supports, and let
\(P,Q\) denote the positive and negative supports. The downsets
\(D\setminus\operatorname{Up}(Q)\) and
\(D\setminus\operatorname{Up}(P)\) differ by exactly \(d\).
Indeed, any point of \(D\) above a point of \(P\) lies below
another point of \(P\), by cross-sign incomparability, and hence
belongs to \(P\) by order convexity; similarly for \(Q\).

At the dyadic character with kernel order \(L\), normalize the
partial difference and sum by \(L\). The difference is \(d\),
and call the sum \(g\). Homometry gives disjoint Fourier supports
for \(d\) and \(g\). Above a size-\(L\) dispute one side stays
empty and the other is empty or has size \(L\): a nonzero smaller
size would be a smaller dispute. Thus \(g=|d|\) on the upward
closure of the support. This is a terminal interpolant, contrary to
the hypothesis.
\end{proof}

\paragraph{Computed result (all prime values).}
The exact symbolic certificate described below excludes every nonzero
equal-population downset difference for \(K=p^2q^2rs\), with
\(p,q,r,s\) distinct odd primes. Together with
Lemma~\ref{lem:terminal-implies-rigidity}, it establishes rigidity of
\(2^ep^2q^2rs\) at every dyadic height, subject to correctness of the
finite exhaustive implementation. We label this result \textsc{Computed}
to keep that implementation dependency visible. It is an unbounded-prime
symbolic certificate, not a test of finitely many numerical prime tuples.

The grid has 36 cells. Enumerate its downsets in a linear extension and
retain all two-sided differences up to overall sign. Each of the six
placements of the two squared axes among increasing prime values gives
11,194,933 distinct signed patterns. For a 36-cell downset with bits
\(b_i\), the key \(\sum_i b_i3^i\) encodes its differences without
loss by uniqueness of balanced ternary. Taking the absolute difference
fixes the orientation; adding \((3^{36}-1)/2\) converts signed digits
to ordinary ternary digits. This arithmetic fits signed 64-bit integers.

For increasing primes use
\(p_i=3+2i+g_0+\cdots+g_i\), with zero-based indices and
nonnegative even gaps \(g_i\). The signed shell transforms are integer
polynomials in these gaps. A polynomial with a nonzero constant and all
coefficients of that weak sign cannot vanish; a polynomial with odd
constant cannot vanish at even gaps. These tests leave a superset of
the possible zero conductors. Population balance requires a possible
zero at the trivial character.

Two further tests exclude possible terminal interpolants. First, at a
disputed shell \(v\), the difference of kernel rows \(v_i\) and
\(v_i-1\), divided by \(p_i^{v_i}\), has coefficients one at
physical order \(v_i\), minus one at \(v_i+1\), and zero elsewhere.
Take their product over positive coordinates, and fix arbitrary single
kernel rows in zero coordinates. This functional is supported in the
prescribed upward region. Whenever its value on \(|d|\) is provably
nonzero, at least one conductor of its row cube must be a possible
Fourier zero. Second, compute the integer-scaled null vectors from
Lemma~\ref{lem:terminal-incidence}. If an identity has exactly one
nonzero term with least exponent on an axis, reduction modulo that
axis prime is impossible when its nonzero integer coefficient has
absolute value smaller than the least possible prime on that axis.

The remaining polynomials are factored over the rationals. A factor is
discarded as nonvanishing only when reduction modulo a selected prime
is a nonzero power of two times a product of known units: other prime
values and smaller prime predecessors. For reduction modulo the largest
prime one may also use smaller-prime successors, which are strictly
smaller than the largest odd prime. Removing such impossible Fourier
zeros permits a further pass of the same necessary tests.

The successive exclusion counts are as follows; omitted categories have
count zero. The exponents are in increasing prime order.
\begin{center}
\small
\begin{tabular}{lrrrrrr}
\hline
Shape & Root survivors & Cube & Incidence & Root factor & Row factor & Largest prime\\
\hline
\(1122\)&2915359&2914533&817&9&0&0\\
\(1212\)&3743527&3742267&1248&12&0&0\\
\(1221\)&4326542&4324839&1690&13&0&0\\
\(2112\)&3772758&3771046&1688&24&0&0\\
\(2121\)&4305351&4303010&2312&28&1&0\\
\(2211\)&5166787&5163679&3072&31&4&1\\
\hline
\end{tabular}
\end{center}
For each row, the exclusions sum to the root survivors. The final
pattern in shape \(2211\), with increasing primes \(p<q<r<s\),
would require
\[
p^2q^2(r-s)-p^2rs+p^2s+q^2rs-q^2r=0.
\]
Modulo \(s\), this forces
\(s\mid q^2r(p-1)(p+1)\), impossible because every displayed
prime or linear factor is positive and smaller than \(s\).

All finite-array coefficient sums are bounded by
\(36\cdot13^6<2^{31}\). Nullspaces and factorizations use exact
rational arithmetic. The supplementary verifier and its dependencies
reproduce all six rows, as part of a 24-grid replay, with the command
\begin{center}
\small\texttt{python3 experiments/level1\_grid\_gap\_certificate.py}
\end{center}
The referee-support document records the pinned source version, runtime
versions, completed checks, and exact output. No implication from a
passing terminal interpolant to an actual mate is used.
